\chapter{Related Work}
%%%%%%%%%%%%%%%%%%%%%%

This chapter surveys prior work across five domains that the proposed system draws on: adaptive threshold methods for HID devices, hyperparameter and threshold optimization, temporal sequence labeling from biosignals, programmatic labeling and weak supervision, and dilated convolutional architectures for sequence modeling.

\section{Adaptive Threshold Methods for HID Devices}

The problem of optimizing actuation parameters for human input devices has a long history in HCI, though prior approaches have relied on fixed heuristics rather than per-user learned models.

Kim et al.~\cite{impact2018} introduced \emph{impact activation}, an alternative actuation scheme for mouse buttons in which a click is registered at the first contact force peak following a rapid finger descent, rather than at a fixed displacement threshold. Their evaluation showed faster reaction times and fewer double-click errors under time pressure. Importantly, impact activation still relies on a single fixed heuristic (the first force peak), making it unsuitable for users with atypical contact dynamics, and it does not adapt to within-session fatigue. This thesis learns thresholds from the user's own recorded telemetry rather than applying a globally fixed policy.

Trewin~\cite{dynamic_keyboard} proposed the \emph{Dynamic Keyboard}, which automatically adjusts keyboard response parameters (minimum press duration, maximum bounce interval) to accommodate users with motor impairments, by observing an initial calibration sequence and adjusting parameters using rule-based adaptation. The Dynamic Keyboard targets binary on/off switch events and adjusts temporal debounce windows; its adaptation rules are domain-expert heuristics rather than models trained from user data, which was a natural design choice given that sub-millimetre analog depth sensing was not commercially available at the time of publication. The proposed system is complementary in motivation---both aim to reduce unintended actuations---but operates at a different signal level (continuous depth at 1\,kHz vs.\ binary event timing) and uses a data-driven model rather than hand-coded rules.

Current commercial implementations~\cite{wooting_rapid_trigger, logitech_superstrike_2026} expose RT and AP as user-adjusted controls without closed-loop personalization, relying on subjective trial-and-error rather than recorded telemetry. The firmware-faithful threshold optimizer in this thesis treats RT/AP tuning as an explicitly optimized per-user objective, closing the loop that these products leave open.

\section{Hyperparameter and Threshold Optimization}

The broader question of how to search continuous parameter spaces with limited evaluation budgets has been extensively studied in the context of machine learning hyperparameter optimization (HPO). Bergstra et al.~\cite{bergstra_tpe_2011} introduced the Tree-structured Parzen Estimator (TPE), a Bayesian approach that models the density of good and bad hyperparameter configurations and selects candidates likely to improve the current optimum. Optuna~\cite{optuna_akiba_2019} operationalizes TPE in a modern framework with dynamic search space definition, pruning of unpromising trials, and distributed execution. Compared to population-based strategies such as CMA-ES~\cite{hansen_cma_2016}, TPE is more sample-efficient in low- to medium-dimensional spaces and is better suited to workloads where each trial is expensive---for instance, when evaluating a candidate AP/RT pair requires full-length replay of telemetry through a firmware finite-state machine (FSM) simulator.

The threshold optimization problem here differs from standard HPO in a few concrete ways: the search space is two-dimensional (AP and RT), each trial evaluation must replay telemetry through a faithful software replica of the peripheral's debounce and state-transition logic, and the objective is a behavioral count (false actuations and missed clicks) rather than a differentiable loss. Optuna's sampling algorithm is used directly; what this thesis adds is the domain-specific trial evaluator that sits inside the study.

\section{Temporal Sequence Labeling in Biosignal Processing}

Detecting click intent from continuous depth telemetry is an instance of temporal sequence labeling, a problem that has received substantial attention in the context of richer biosignals.

Surface electromyography (sEMG) is the closest related modality. Wang et al.~\cite{reactemg} presented ReactEMG, which uses masked autoencoder pretraining on sEMG to achieve stable, low-latency intent detection from eight-channel surface electrode recordings. Antuvan and Masia~\cite{lda_emg} applied Linear Discriminant Analysis to simultaneous classification of wrist and finger movements from sEMG, demonstrating real-time throughput on embedded hardware. Both systems operate on multi-channel electrode signals that encode rich neuromuscular information before mechanical actuation occurs; their signal-to-noise advantage over single-sensor depth telemetry is substantial. The proposed system, by contrast, uses only the depth signal already present in the peripheral itself, requiring no additional sensing hardware. This imposes a harder classification problem---depth telemetry cannot observe pre-actuation motor intent---but removes the deployment barrier of wearable electrodes.

Keystroke dynamics literature~\cite{typeformer2022, continuous_keystroke_mdpi} models typing patterns for biometric authentication from binary key event timestamps. TypeFormer~\cite{typeformer2022} applies a Transformer architecture to sequences of inter-key timings and achieves high verification accuracy on mobile devices. These methods treat the keystroke as an atomic event and model the temporal pattern between events; they do not model sub-millisecond intra-click dynamics. The proposed system operates at a finer granularity---frame-level classification at 1\,kHz---which requires a different feature engineering approach and exposes a different class of failure modes (contact-bounce artefacts, quantization noise) absent from binary event streams.

\section{Weak Supervision and Algorithmic Bootstrapping}

The geometric auto-labeler introduced in Section~\ref{sec:labeling} fundamentally differs from programmatic weak supervision frameworks such as Snorkel~\cite{ratner_snorkel_2020}. Snorkel defines \emph{labeling functions}---heuristic rules, distant supervision signals, or domain-expert logic---that each assign noisy labels to unlabeled examples, and combines them via a generative model that learns the accuracy and correlations of each labeling function. The combined label distribution is then used as the training signal for a downstream discriminative model.

Instead of deploying the geometric auto-labeler as an automated weak supervision signal for model training, the proposed system employs it strictly as an algorithmic bootstrapping assistant for human-in-the-loop annotation. As a result, the proposed system acts as a strongly supervised learning pipeline accelerated by algorithmic pre-labeling, rather than a weakly supervised system relying on noisy heuristics. To account for natural human inconsistency in selecting the exact millisecond of an onset during this manual process, the pipeline also supports asymmetric soft labels (described in Section~\ref{sec:modeling}) that would encode physiological directional uncertainty into the verified targets rather than treating boundary annotations as rigid, perfect temporal observations.

\section{Dilated Convolutions for Sequence Modeling}

The OfflineFCN architecture adapts the dilated causal convolution structure introduced by van den Oord et al.\ in WaveNet~\cite{van_den_oord_wavenet_2016} for raw audio generation. WaveNet stacks dilated causal convolutional layers with exponentially increasing dilation rates to achieve a large receptive field with manageable parameter counts. The dilation schedule $d_i = 2^i$ allows the stack to cover $\sum 2^i$ steps of context with $O(\log N)$ layers rather than the $O(N)$ layers that a fixed-dilation architecture would require.

Bai et al.~\cite{bai2018tcn} conducted a systematic empirical comparison of Temporal Convolutional Networks (TCNs) and recurrent architectures (LSTMs, GRUs) across a range of sequence modeling benchmarks, demonstrating that dilated causal CNNs match or exceed recurrent models in accuracy while training faster and avoiding gradient vanishing. Their work established the TCN as a practical default for 1-D sequence classification, and the OfflineFCN design is motivated by the same receptive-field arguments they formalize. Ismail Fawaz et al.\ further benchmarked deep learning classifiers for time-series classification and found that FCN-based and ResNet-based architectures consistently outperform LSTM-based models on sensor data~\cite{ismail_fawaz_ts_2019}, supporting the choice to build around convolutional blocks rather than recurrent ones.

The OfflineFCN departs from the standard causal TCN in one critical dimension: it uses symmetric (non-causal) padding rather than causal (left-only) padding. This doubles the effective receptive field per layer for the same kernel size, since context is drawn from both directions---enabling a large symmetric receptive field, described in Section~\ref{sec:modeling}, to be achieved with only 8 blocks. The trade-off is that the resulting model cannot be deployed for streaming inference, since each output depends on future samples. This is appropriate for the offline analysis mode but excludes the model from the on-device path described as future work.


%%%%%%%%%%%%%%%%%%%%
